% Complete reviewed Unit 14; editorial master: chapters/ch14/answer-key.md
% Markdown SHA-256: c109f905cd5a08fe5e734426ea61202b8ae7888845ae80682bc31057bc17886b
% One-use initial assembly followed by complete content and page review.
\newcommand{\runningtitle}{Unit 14. Answer key}
% Layout inherited from the reviewed previous course; set a running title per unit.
% This is a small hand-maintained TeX layout, not a conversion pipeline.
\RequirePackage{fix-cm}
\documentclass[11pt,twoside]{article}
\usepackage{fontspec}
\setmainfont{Cambria}[
  Renderer=HarfBuzz,
  BoldFont=cambriab.ttf,
  ItalicFont=cambriai.ttf,
  BoldItalicFont=cambriaz.ttf
]
\setmonofont{lmmono10-regular.otf}[BoldFont=lmmonolt10-bold.otf]
\usepackage[paperwidth=165mm,paperheight=241mm,inner=21mm,outer=19mm,
  top=22mm,bottom=20mm,headheight=14pt,headsep=7mm,footskip=10mm]{geometry}
\usepackage{xcolor,array,booktabs,longtable,tabularx}
\usepackage{enumitem,fancyhdr}
\usepackage[unicode,hidelinks]{hyperref}
\definecolor{ink}{HTML}{000000}
\definecolor{muted}{HTML}{4C5558}
\setlength{\parindent}{0pt}
\setlength{\parskip}{3.5pt}
\AtBeginDocument{\fontsize{11}{14.3}\selectfont}
\setlength{\tabcolsep}{4pt}
\renewcommand{\arraystretch}{1.10}
\setlist[itemize]{leftmargin=1.3em,label=\textbullet,itemsep=2pt,topsep=3pt,parsep=0pt}
\setlist[enumerate]{leftmargin=1.8em,itemsep=3pt,topsep=4pt,parsep=0pt}
\widowpenalty=10000
\clubpenalty=10000
\emergencystretch=1em
\raggedbottom
\pagestyle{fancy}
\fancyhf{}
\providecommand{\runningtitle}{Old English}
\fancyhead[LE,RO]{\small\itshape \runningtitle}
\fancyhead[LO,RE]{\small Old English}
\fancyfoot[LE,RO]{\small\thepage}
\renewcommand{\headrulewidth}{0.3pt}
\newcommand{\booktitle}[2]{
  {\fontsize{9}{11}\selectfont\color{muted}#1\par}
  \vspace{3pt}
  {\fontsize{20}{24}\selectfont\color{ink}#2\par}
  \vspace{6pt}
}
\newcommand{\source}[1]{\par{\footnotesize\color{muted}#1\par}}
\newcommand{\readingtext}[1]{
  \par\addvspace{5pt}
  \begingroup\leftskip=1.5em
    \fontsize{12}{16}\selectfont #1\par
  \endgroup\addvspace{5pt}
}
\makeatletter
\renewcommand{\section}{\@startsection{section}{1}{\z@}
  {-2.4ex \@plus -0.6ex \@minus -0.2ex}{0.8ex \@plus 0.2ex}
  {\normalfont\fontsize{12}{15}\selectfont\bfseries\color{ink}}}
\renewcommand{\subsection}{\@startsection{subsection}{2}{\z@}
  {-1.8ex \@plus -0.5ex \@minus -0.2ex}{0.5ex \@plus 0.2ex}
  {\normalfont\fontsize{11}{14}\selectfont\bfseries\color{ink}}}
\makeatother

\hyphenpenalty=10000
\exhyphenpenalty=10000
\emergencystretch=3em
\newcommand{\itempara}[2]{\par\addvspace{3pt}\noindent\hangindent=1.8em\hangafter=1\makebox[1.8em][l]{#1.}#2\par}
\newcommand{\reserve}[1]{\par\penalty-100\begingroup\dimen0=\pagegoal\advance\dimen0 by-\pagetotal\ifdim\dimen0<#1\newpage\fi\endgroup}
\begin{document}
\addfontfeatures{Ligatures=TeXOff}
\raggedright
\booktitle{UNIT 14}{Answer key}

% mdparagraph:P001:begin
This key covers all 44 tasks: 23 original exercises and the two complete Cook sets. The models, translations and explanations are new teaching material. The four translated teaching examples are not additional exercises. Equivalent English wording and defensible Old English alternatives are acceptable. Judge a stated vocabulary, form or word-order target separately from grammatical validity. [1]
% mdparagraph:P001:end

% mdparagraph:P002:begin
In form tasks, give every labelled cell even when several cells have the same spelling. In composition, preserve the participants, case relations, number, mood and intended meaning. A prefix belongs to a particular lexical item; it does not supply a universal English meaning or make its verb interchangeable with the unprefixed verb. Supplied words and phrases remain available throughout the tasks and are not tests of unaided recall.
% mdparagraph:P002:end

\clearpage
% mdparagraph:P003:begin
\section*{F. Complete forms}
% mdparagraph:P003:end

\reserve{5\baselineskip}
\reserve{8\baselineskip}
% mdparagraph:P004:begin
\textbf{F01. Wǣdla and ēage.}
% mdparagraph:P004:end

% item: NEW-U14-F01
% key: KEY-NEW-U14-F01
\reserve{11\baselineskip}
% mdparagraph:P005:begin
\textbf{Wǣdla}, poor man, is weak masculine.
% mdparagraph:P005:end

% mdtable:T01:begin
\par\begingroup
\fontsize{10.5}{13.2}\selectfont
\setlength{\LTleft}{0pt}\setlength{\LTright}{\fill}
\noindent
\begin{tabular}{@{}>{\raggedright\arraybackslash}p{39.46mm}>{\raggedright\arraybackslash}p{39.46mm}>{\raggedright\arraybackslash}p{39.46mm}@{}}
\toprule
\textbf{Case} & \textbf{Singular} & \textbf{Plural} \\
\midrule
Nominative & wǣdla & wǣdlan \\
Genitive & wǣdlan & wǣdlena \\
Dative & wǣdlan & wǣdlum \\
Accusative & wǣdlan & wǣdlan \\
\bottomrule
\end{tabular}
\par\endgroup

% mdtable:T01:end
\reserve{11\baselineskip}
% mdparagraph:P006:begin
\textbf{Ēage}, eye, is weak neuter.
% mdparagraph:P006:end

% mdtable:T02:begin
\par\begingroup
\fontsize{10.5}{13.2}\selectfont
\setlength{\LTleft}{0pt}\setlength{\LTright}{\fill}
\noindent
\begin{tabular}{@{}>{\raggedright\arraybackslash}p{39.46mm}>{\raggedright\arraybackslash}p{39.46mm}>{\raggedright\arraybackslash}p{39.46mm}@{}}
\toprule
\textbf{Case} & \textbf{Singular} & \textbf{Plural} \\
\midrule
Nominative & ēage & ēagan \\
Genitive & ēagan & ēagena \\
Dative & ēagan & ēagum \\
Accusative & ēage & ēagan \\
\bottomrule
\end{tabular}
\par\endgroup

% mdtable:T02:end
% mdparagraph:P007:begin
The masculine distinguishes nominative singular \textbf{-a} from accusative singular \textbf{-an}. The neuter has \textbf{-e} in both singular nominative and accusative. In wǣdla, \textbf{-an} can mark singular genitive, dative or accusative, or plural nominative or accusative. In ēage, it can mark singular genitive or dative, or plural nominative or accusative. Context supplies the distinction; -an alone does not mean “plural” or “object”. The documented genitive plural \textbf{ēagna} is also valid. Do not extend either the masculine -a or its accusative -an into the neuter's singular direct cases. [2]
% mdparagraph:P007:end

\reserve{11\baselineskip}
\reserve{9\baselineskip}
\par\begin{minipage}{\linewidth}\raggedright
% mdparagraph:P008:begin
\textbf{F02. Compound nouns.}
% mdparagraph:P008:end

% item: NEW-U14-F02
% key: KEY-NEW-U14-F02
% mdtable:T03:begin
\par\begingroup
\fontsize{10.5}{13.2}\selectfont
\setlength{\LTleft}{0pt}\setlength{\LTright}{\fill}
\noindent
\begin{tabular}{@{}>{\raggedright\arraybackslash}p{39.46mm}>{\raggedright\arraybackslash}p{39.46mm}>{\raggedright\arraybackslash}p{39.46mm}@{}}
\toprule
\textbf{Case} & \textbf{Singular} & \textbf{Plural} \\
\midrule
Nominative & goldhord & goldhord \\
Genitive & goldhordes & goldhorda \\
Dative & goldhorde & goldhordum \\
Accusative & goldhord & goldhord \\
\bottomrule
\end{tabular}
\par\endgroup

% mdtable:T03:end
% mdparagraph:P009:begin
The other two dative singulars are \textbf{wyrtgemange} and \textbf{gōdwębbe}. Attach the ending to the whole compound: \textbf{goldhord-e, wyrtgemang-e, gōdwębb-e}. The first element does not receive a separate case ending inside the compound. The heavy neuter goldhord has no -u in its nominative/\allowbreak{}accusative plural; context distinguishes those forms from the singular. Preserve the long ō in gōdwębb. These are inflected forms of the same nouns, not new lexical derivatives. [2]
% mdparagraph:P009:end

\end{minipage}\par\addvspace{4pt}

\par\begin{minipage}{\linewidth}\raggedright
% mdparagraph:P010:begin
\textbf{F03. Liccian.}
% mdparagraph:P010:end

% item: NEW-U14-F03
% key: KEY-NEW-U14-F03
% mdtable:T04:begin
\par\begingroup
\fontsize{10.5}{13.2}\selectfont
\setlength{\LTleft}{0pt}\setlength{\LTright}{\fill}
\noindent
\begin{tabular}{@{}>{\raggedright\arraybackslash}p{33.63mm}>{\raggedright\arraybackslash}p{41.02mm}>{\raggedright\arraybackslash}p{43.73mm}@{}}
\toprule
\textbf{Person /\allowbreak{} number} & \textbf{Present indicative} & \textbf{Preterite indicative} \\
\midrule
1 singular & liccie & liccode \\
2 singular & liccast & liccodest \\
3 singular & liccað & liccode \\
1 plural & licciað & liccodon \\
2 plural & licciað & liccodon \\
3 plural & licciað & liccodon \\
\bottomrule
\end{tabular}
\par\endgroup

% mdtable:T04:end
% mdparagraph:P011:begin
The selected weak-II pattern has \textbf{licci-} in the first singular and all plural present forms, but \textbf{licc-} before second-singular -ast and third-singular -að. The preterite has \textbf{liccod-} before -e, -est and -on. Both c letters remain. A plural subject needs licciað or liccodon, even though the three persons within either plural tense coincide. Do not form a strong preterite by changing the root vowel. [3]
% mdparagraph:P011:end

\end{minipage}\par\addvspace{4pt}

\reserve{8\baselineskip}
% mdparagraph:P012:begin
\textbf{F04. Selected indicative forms.}
% mdparagraph:P012:end

% item: NEW-U14-F04
% key: KEY-NEW-U14-F04
% mdtable:T05:begin
\par\begingroup
\fontsize{10.5}{13.2}\selectfont
\setlength{\LTleft}{0pt}\setlength{\LTright}{\fill}
\noindent
\begin{tabular}{@{}>{\raggedright\arraybackslash}p{78.74mm}>{\raggedright\arraybackslash}p{42.45mm}@{}}
\toprule
\textbf{Verb and requested cell} & \textbf{Form} \\
\midrule
gewęndan, present 2 singular & gewęntst \\
gelīefan, present 3 singular & gelīefð \\
forgiefan, present 3 singular & forgiefð \\
āhrēosan, preterite 3 plural & āhruron \\
gebiddan, present 3 singular & gebitt \\
ætīewan, preterite 3 plural & ætīewdon \\
āhębban, present 3 singular & āhęfð \\
onfōn, preterite 2 singular & onfēnge \\
\bottomrule
\end{tabular}
\par\endgroup

% mdtable:T05:end
% mdparagraph:P013:begin
\textbf{Āhruron} uses the strong-II preterite plural stem \textbf{āhrur-}, including its s/\allowbreak{}r alternation. It is not made by adding -on to the present stem āhrēos-. \textbf{Onfēnge} uses \textbf{onfēng-} plus second-singular -e. A strong preterite second singular belongs with the preterite plural stem, even though its number is singular. The attested prefixes remain part of both verbs. [3]
% mdparagraph:P013:end

% mdparagraph:P014:begin
Keep the diagnostic changes in the other cells: gewęnd- gives \textbf{gewęntst}; gebidd- gives \textbf{gebitt}; āhębb- gives \textbf{āhęfð}. \textbf{Ætīewdon} belongs to the weak-I verb “show”; the earlier weak-II verb “appear” has a different system. Long īe in gelīefð does not license adding a length mark to forgiefð.
% mdparagraph:P014:end

\reserve{13\baselineskip}
\reserve{11\baselineskip}
\par\begin{minipage}{\linewidth}\raggedright
% mdparagraph:P015:begin
\textbf{F05. Subjunctives.}
% mdparagraph:P015:end

% item: NEW-U14-F05
% key: KEY-NEW-U14-F05
% mdtable:T06:begin
\par\begingroup
\fontsize{10.5}{13.2}\selectfont
\setlength{\LTleft}{0pt}\setlength{\LTright}{\fill}
\noindent
\begin{tabular}{@{}>{\raggedright\arraybackslash}p{32.0mm}>{\raggedright\arraybackslash}p{41.9mm}>{\raggedright\arraybackslash}p{44.48mm}@{}}
\toprule
\textbf{Person /\allowbreak{} number} & \textbf{Present subjunctive} & \textbf{Preterite subjunctive} \\
\midrule
1 singular & gebidde & gebǣde \\
2 singular & gebidde & gebǣde \\
3 singular & gebidde & gebǣde \\
1 plural & gebidden & gebǣden \\
2 plural & gebidden & gebǣden \\
3 plural & gebidden & gebǣden \\
\bottomrule
\end{tabular}
\par\endgroup

% mdtable:T06:end
% mdtable:T07:begin
\par\begingroup
\fontsize{10.5}{13.2}\selectfont
\setlength{\LTleft}{0pt}\setlength{\LTright}{\fill}
\noindent
\begin{tabular}{@{}>{\raggedright\arraybackslash}p{86.51mm}>{\raggedright\arraybackslash}p{34.67mm}@{}}
\toprule
\textbf{Additional requested cell} & \textbf{Form} \\
\midrule
forweorðan, present subjunctive 3 singular & forweorðe \\
sęndan, present subjunctive 2 singular & sęnde \\
gelīefan, present subjunctive 3 plural & gelīefen \\
ārīsan, present subjunctive 3 singular & ārīse \\
\bottomrule
\end{tabular}
\par\endgroup

% mdtable:T07:end
% mdparagraph:P016:begin
All persons share the singular form, and all persons share the plural form. The subjunctive has no distinct second-singular -st. \textbf{Gebǣde/\allowbreak{}gebǣden} use the strong preterite plural stem with long ǣ; the present keeps \textbf{gebidd-}. The spelling sęnde can also occur elsewhere in the verb's system: its label and clause determine which cell is intended. This task asks for forms; mood choice in a sentence must still follow that sentence's construction. [3]
% mdparagraph:P016:end

\end{minipage}\par\addvspace{4pt}

\par\begin{minipage}{\linewidth}\raggedright
% mdparagraph:P017:begin
\textbf{F06. Participles and commands.}
% mdparagraph:P017:end

% item: NEW-U14-F06
% key: KEY-NEW-U14-F06
% mdtable:T08:begin
\par\begingroup
\fontsize{10.5}{13.2}\selectfont
\setlength{\LTleft}{0pt}\setlength{\LTright}{\fill}
\noindent
\begin{tabular}{@{}>{\raggedright\arraybackslash}p{28.56mm}>{\raggedright\arraybackslash}p{48.41mm}>{\raggedright\arraybackslash}p{41.41mm}@{}}
\toprule
\textbf{Verb} & \textbf{Present participle base} & \textbf{Past participle base} \\
\midrule
forceorfan & forceorfende & forcorfen \\
weorpan & weorpende & worpen \\
gemētan & gemētende & gemētt \\
geglęngan & geglęngende & geglęnged \\
\bottomrule
\end{tabular}
\par\endgroup

% mdtable:T08:end
% mdtable:T09:begin
\par\begingroup
\fontsize{10.5}{13.2}\selectfont
\setlength{\LTleft}{0pt}\setlength{\LTright}{\fill}
\noindent
\begin{tabular}{@{}>{\raggedright\arraybackslash}p{35.44mm}>{\raggedright\arraybackslash}p{41.11mm}>{\raggedright\arraybackslash}p{41.82mm}@{}}
\toprule
\textbf{Verb} & \textbf{Singular command} & \textbf{Plural command} \\
\midrule
gemiltsian & gemiltsa & gemiltsiað \\
tilian & tila & tiliað \\
\bottomrule
\end{tabular}
\par\endgroup

% mdtable:T09:end
% mdparagraph:P018:begin
The present bases use -ende with the present stem; the strong past bases use the fourth stem. Thus forceorf- gives \textbf{forcorfen}, and weorp- gives \textbf{worpen}. The weak bases retain their own dental formation, including \textbf{gemētt}. \textbf{Geworpen} is a documented alternative to bare worpen. Do not add another ge- to gemētt or geglęnged merely because a form is a participle. [3]
% mdparagraph:P018:end

% mdparagraph:P019:begin
The commands have singular -a and plural -iað. Preserve \textbf{ts} in gemiltsian and the single \textbf{l} in tilian. The participle task requests bases, not full attributive adjective declensions or new finite tenses.
% mdparagraph:P019:end

\end{minipage}\par\addvspace{4pt}

\reserve{7\baselineskip}
% mdparagraph:P020:begin
\section*{A. Word families and analysis}
% mdparagraph:P020:end

\reserve{10\baselineskip}
\reserve{8\baselineskip}
\par\begin{minipage}{\linewidth}\raggedright
% mdparagraph:P021:begin
\textbf{A01. Compounds and endings.}
% mdparagraph:P021:end

% item: NEW-U14-A01
% key: KEY-NEW-U14-A01
% mdtable:T10:begin
\par\begingroup
\fontsize{10.5}{13.2}\selectfont
\setlength{\LTleft}{0pt}\setlength{\LTright}{\fill}
\noindent
\begin{tabular}{@{}>{\raggedright\arraybackslash}p{22mm}>{\raggedright\arraybackslash}p{32mm}>{\raggedright\arraybackslash}p{64.37mm}@{}}
\toprule
\textbf{Expression} & \textbf{Components and head} & \textbf{Inflection and meaning} \\
\midrule
goldhord & gold + hord; final head hord & The citation has no overt ending: treasure, a hoard of gold. \\
heofonfȳre & heofon + fȳr; final head fȳr & -e is dative singular: fire from heaven, here destructive heavenly fire/\allowbreak{}lightning. \\
hēaclifum & hēa- + clif; final head clif & -um is dative plural: high cliffs. \\
\bottomrule
\end{tabular}
\par\endgroup

% mdtable:T10:end
% mdparagraph:P022:begin
The first element specifies the final noun's content, source or quality. In hēaclifum, \textbf{hēa-} is the attested combining form related to hēah. The visible relationship does not make dropping h a general rule for new words. Components help explain an attested word; lexical use still determines its exact meaning. The noun citations hord and clif and the two source forms remain supplied. Do not treat -e or -um as part of the lexical compound, or infer a complete new declension from either source token. [4]
% mdparagraph:P022:end

\end{minipage}\par\addvspace{4pt}

\par\begin{minipage}{\linewidth}\raggedright
% mdparagraph:P023:begin
\textbf{A02. Derived nouns.}
% mdparagraph:P023:end

% item: NEW-U14-A02
% key: KEY-NEW-U14-A02
% mdtable:T11:begin
\par\begingroup
\fontsize{10.5}{13.2}\selectfont
\setlength{\LTleft}{0pt}\setlength{\LTright}{\fill}
\noindent
\begin{tabular}{@{}>{\raggedright\arraybackslash}p{29.17mm}>{\raggedright\arraybackslash}p{23.53mm}>{\raggedright\arraybackslash}p{65.68mm}@{}}
\toprule
\textbf{Family} & \textbf{Formation} & \textbf{Parts of speech and meanings} \\
\midrule
leornian /\allowbreak{} leornung & leorn- + -ung & Verb “learn”; feminine noun “learning, study”. \\
gefēra /\allowbreak{} gefērscipe & gefēr- + -scipe & Noun “companion”; masculine noun “companionship”. \\
rihtwīs /\allowbreak{} rihtwīsnes & rihtwīs + -nes & Adjective “righteous”; feminine noun “righteousness”. \\
īdel /\allowbreak{} īdelnes & īdel + -nes & Adjective “empty, void”; feminine noun “idleness”. \\
\bottomrule
\end{tabular}
\par\endgroup

% mdtable:T11:end
% mdparagraph:P024:begin
The weak-noun citation ending -a of gefēra is not retained before -scipe. Each result is a lexical noun with its own gender, sense and inflection. \textbf{Īdelnes} has the taught sense “idleness”; mechanically translating the earlier adjective sense plus an English suffix as “emptiness” misses that lexical choice. “Indolence” is a suitable equivalent. These examples establish attested family relationships, not a method for inventing an unrestricted set of nouns. [4]
% mdparagraph:P024:end

\end{minipage}\par\addvspace{4pt}

\par\begin{minipage}{\linewidth}\raggedright
% mdparagraph:P025:begin
\textbf{A03. Adjectives and adverbs.}
% mdparagraph:P025:end

% item: NEW-U14-A03
% key: KEY-NEW-U14-A03
% mdparagraph:P026:begin
\textbf{Heofon}, heaven/\allowbreak{}sky, is a noun; \textbf{-lic} forms the adjective \textbf{heofonlic}, heavenly. \textbf{Swīð}, strong, is an adjective citation; \textbf{swīðe}, very, is an adverb with -e. \textbf{Sōð}, true, is an adjective citation; \textbf{sōðlice}, truly/\allowbreak{}indeed, is an adverb with -lice. The adverbs do not agree with a noun. Swīðe's selected meaning is “very”, not a mechanically substituted “strong”.
% mdparagraph:P026:end

% mdparagraph:P027:begin
A final e alone proves nothing about word class: it can also be a nominal or adjectival case ending. The whole word, its lexical entry and its function settle the analysis. The supplied swīð and sōð citations do not introduce complete adjective paradigms. [4]
% mdparagraph:P027:end

\end{minipage}\par\addvspace{4pt}

\par\begin{minipage}{\linewidth}\raggedright
% mdparagraph:P028:begin
\textbf{A04. Prefixes and whole-word meanings.}
% mdparagraph:P028:end

% item: NEW-U14-A04
% key: KEY-NEW-U14-A04
% mdparagraph:P029:begin
\textbf{Un-} marks opposition in \textbf{rihtwīsnes /\allowbreak{} unrihtwīsnes}, righteousness/\allowbreak{}unrighteousness. \textbf{Tō-} contributes the attested “asunder” relationship in \textbf{berstan /\allowbreak{} tōberstan}, burst/\allowbreak{}break up or go to pieces. \textbf{For-} belongs to \textbf{forweorðan}, perish, beside \textbf{weorðan}, become. “Perish” cannot be recovered by translating the prefix as the English preposition “for”.
% mdparagraph:P029:end

% mdparagraph:P030:begin
The whole word's meaning and government must be learned. Prefixes have different functions in different words; these pairs neither prove that another imagined compound exists nor permit replacing every prefixed verb with its base. The four verb citations are supplied for analysis. No additional syntactic infinitive use is being tested. [4]
% mdparagraph:P030:end

\end{minipage}\par\addvspace{4pt}

\par\begin{minipage}{\linewidth}\raggedright
% mdparagraph:P031:begin
\textbf{A05. Derivation and inflection.}
% mdparagraph:P031:end

% item: NEW-U14-A05
% key: KEY-NEW-U14-A05
% mdparagraph:P032:begin
\textbf{Rihtwīs → rihtwīsnes} is derivation: -nes forms the noun “righteousness” from the adjective “righteous”. \textbf{Rihtwīsnes → rihtwīsnesse} is inflection: the noun uses its oblique stem \textbf{rihtwīsness-} with -e. It remains the same lexical noun.
% mdparagraph:P032:end

% mdparagraph:P033:begin
In \textbf{Ic lufie rihtwīsnesse}, “I love righteousness”, rihtwīsnesse is feminine accusative singular, the object of lufie. In isolation that spelling can also be singular genitive or dative. Neither the final -e nor the abstract meaning identifies its case without the clause. Confusing the derivational -nes with the following inflection misses the two different operations. [2, 4]
% mdparagraph:P033:end

\end{minipage}\par\addvspace{4pt}

\par\begin{minipage}{\linewidth}\raggedright
% mdparagraph:P034:begin
\textbf{A06. A lexical noun and a substantival adjective.}
% mdparagraph:P034:end

% item: NEW-U14-A06
% key: KEY-NEW-U14-A06
% mdparagraph:P035:begin
\textbf{Hē onfēng gōdum:} “He received goods/\allowbreak{}possessions.” Gōdum is dative plural of the assigned neuter noun gōd. Onfōn governs the dative thing received; an English direct object does not establish an Old English accusative.
% mdparagraph:P035:end

% mdparagraph:P036:begin
\textbf{Ic geseah ðā rihtwīsan:} “I saw the righteous people.” In the chapter's selected generic masculine people interpretation, rihtwīsan is a weak accusative plural adjective used substantivally. It denotes the people seen and agrees with the article in case, number and gender. The article supplies the definite expression, and the adjective retains its weak inflection. The plural endings do not distinguish gender. In isolation, the same string could also mean “I saw the righteous woman”, with feminine accusative singular; the chapter's explicit people gloss selects the plural reading for this task.
% mdparagraph:P036:end

% mdparagraph:P037:begin
An absent following noun does not by itself decide lexical category. Gōd here is the separately evidenced lexical noun “goods”; rihtwīs has the previously taught adjectival personal use. The same surface ending can occur in different word classes and cases. [5]
% mdparagraph:P037:end

\end{minipage}\par\addvspace{4pt}

\reserve{7\baselineskip}
% mdparagraph:P038:begin
\section*{C. Controlled composition}
% mdparagraph:P038:end

% mdparagraph:P039:begin
The stated vocabulary is part of each target. Word order can vary where the case relations and reference remain clear. A defensible synonym may express the sentence correctly while missing the requested word.
% mdparagraph:P039:end

\reserve{5\baselineskip}
\par\begin{minipage}{\linewidth}\raggedright
% mdparagraph:P040:begin
\textbf{C01. Poor men, dogs and wounds.}
% mdparagraph:P040:end

% item: NEW-U14-C01
% key: KEY-NEW-U14-C01
% mdparagraph:P041:begin
\textbf{Ðā wǣdlan gesāwon ðā hundas and ðā wunda.}
% mdparagraph:P041:end

% mdparagraph:P042:begin
The poor men are the nominative plural subject; the dogs and wounds are coordinated accusative plural objects. Gesāwon must be plural. Wǣdlan and hundas have matching nominative/\allowbreak{}accusative plural spellings, but their functions differ here. Wunda belongs to the feminine noun wund; masculine wundas is not its selected plural. Reordering the objects is possible if the same two objects are preserved. [2, 6]
% mdparagraph:P042:end

\end{minipage}\par\addvspace{4pt}

\par\begin{minipage}{\linewidth}\raggedright
% mdparagraph:P043:begin
\textbf{C02. Leaving and giving.}
% mdparagraph:P043:end

% item: NEW-U14-C02
% key: KEY-NEW-U14-C02
% mdparagraph:P044:begin
\textbf{Ic forlǣte ēow sibbe, and ic forgiefe ðǣm ðearfan ðæt goldhord.}
% mdparagraph:P044:end

% mdparagraph:P045:begin
“I leave you peace, and I give the poor man the treasure.” Both recipients are dative: plural ēow and singular ðǣm ðearfan. Sibbe and ðæt goldhord are accusative things transferred. Forlǣtan here has the newly taught leave-peace frame. Fronting sibbe or another complete constituent can be valid; changing ēow to singular ðē misses the number target. The recipient's weak-noun -an does not make it accusative. [5]
% mdparagraph:P045:end

\end{minipage}\par\addvspace{4pt}

\par\begin{minipage}{\linewidth}\raggedright
% mdparagraph:P046:begin
\textbf{C03. Belief.}
% mdparagraph:P046:end

% item: NEW-U14-C03
% key: KEY-NEW-U14-C03
% mdparagraph:P047:begin
\textbf{Ðā wǣdlan gelīefað on God.}
% mdparagraph:P047:end

% mdparagraph:P048:begin
The plural poor men believe in God. Gelīefað agrees with the plural subject; \textbf{on + accusative} belongs to this belief frame. God has the accusative spelling used here. Do not explain this case as physical motion or substitute dative Gode on the basis of English “in”. [5]
% mdparagraph:P048:end

\end{minipage}\par\addvspace{4pt}

\par\begin{minipage}{\linewidth}\raggedright
% mdparagraph:P049:begin
\textbf{C04. Two commands.}
% mdparagraph:P049:end

% item: NEW-U14-C04
% key: KEY-NEW-U14-C04
% mdparagraph:P050:begin
\textbf{Gemiltsa ðǣm wǣdlan, and sęnd him ðæt gold.}
% mdparagraph:P050:end

% mdparagraph:P051:begin
“Have mercy on the poor man, and send him the gold.” Both commands address one person. Gemiltsian governs the dative person; him is the dative recipient of sęnd, referring to the same poor man. The gold is accusative. Repeating ðǣm wǣdlan instead of him can make the reference explicit, but plural gemiltsiað or sęndað misses the singular-command target. [5]
% mdparagraph:P051:end

\end{minipage}\par\addvspace{4pt}

\reserve{7\baselineskip}
% mdparagraph:P052:begin
\section*{Cook VIII.26}
% mdparagraph:P052:end

% mdparagraph:P053:begin
The original prompts, including their parenthetic controls, remain in the learner chapter. These models are new compositions. The local helps remain available. The specified final-verb arrangements are exercise targets; ordinary alternative orders must be judged separately. [6]
% mdparagraph:P053:end

\reserve{5\baselineskip}
\par\begin{minipage}{\linewidth}\raggedright
% mdparagraph:P054:begin
\textbf{VIII.26.1.}
% mdparagraph:P054:end

% item: EX-VIII-26-01
% key: KEY-EX-VIII-26-01
% mdparagraph:P055:begin
Ðū eart dūst, and ðū gewęntst tō dūste.
% mdparagraph:P055:end

% mdparagraph:P056:begin
The same second-person subject is dust and returns to dust. \textbf{Dūst} is a nominative predicate; \textbf{dūste} is dative after tō. \textbf{Gewęntst} has the selected dental contraction. Repeating ðū is clear but not obligatory. Do not give the two occurrences of the noun the same ending merely because English repeats “dust”.
% mdparagraph:P056:end

\end{minipage}\par\addvspace{4pt}

\par\begin{minipage}{\linewidth}\raggedright
% mdparagraph:P057:begin
\textbf{VIII.26.2.}
% mdparagraph:P057:end

% item: EX-VIII-26-02
% key: KEY-EX-VIII-26-02
% mdparagraph:P058:begin
Swā swīðe lufode God middangeard, ðæt hē sealde his āncęnnedan Sunu, ðæt nān ne forweorðe ðe on hine gelīefð, ac hæbbe ðæt ēce līf.
% mdparagraph:P058:end

% mdparagraph:P059:begin
God's love results in his giving the Son; the intended outcome of that gift is that the believer should not perish but have eternal life. \textbf{Sealde}, from the specified sęllan, states the giving. \textbf{Forweorðe} and \textbf{hæbbe} are coordinated purpose subjunctives; no separate “may” verb is needed. The article is omitted before \textbf{middangeard} as requested. The supplied \textbf{nān … ne} expresses one negation. Bare relative \textbf{ðe}, Cook §87.d, introduces the believer's defining action; \textbf{gelīefð on hine} uses on with accusative. [6]
% mdparagraph:P059:end

% mdparagraph:P060:begin
The supplied āncęnnedan and the earlier Son expression preserve the intended referent. A different defensible word order may retain the meaning. An inflected relative followed by ðe may be valid elsewhere, but does not satisfy this bare-relative control. Indicatives for both purpose verbs would change the stated construction and blur the distinction between the gift and its intended result.
% mdparagraph:P060:end

\end{minipage}\par\addvspace{4pt}

\par\begin{minipage}{\linewidth}\raggedright
% mdparagraph:P061:begin
\textbf{VIII.26.3.}
% mdparagraph:P061:end

% item: EX-VIII-26-03
% key: KEY-EX-VIII-26-03
% mdparagraph:P062:begin
Se Hǣlend cwæð tō him, Ic eom weg and sōðfæstnes and līf; nān ne cymð tō ðǣm Fæder būtan ðurh mē.
% mdparagraph:P062:end

% mdparagraph:P063:begin
The direct quotation retains the speaker's \textbf{ic … mē}. \textbf{Weg, sōðfæstnes, līf} are bare nominative predicates identifying him. The first tō phrase gives the person addressed; \textbf{tō ðǣm Fæder} gives the destination. \textbf{Ðurh mē} has accusative mē, while the supplied būtan means “except”. \textbf{Nān ne cymð} is one negation, not an assertion that someone does come. The selected se Hǣlend renders Jesus here; a defensible established personal name could also identify him. Do not add an article to each predicate solely from a familiar modern translation.
% mdparagraph:P063:end

\end{minipage}\par\addvspace{4pt}

\par\begin{minipage}{\linewidth}\raggedright
% mdparagraph:P064:begin
\textbf{VIII.26.4.}
% mdparagraph:P064:end

% item: EX-VIII-26-04
% key: KEY-EX-VIII-26-04
% mdparagraph:P065:begin
Sibbe ic forlǣte ēow; mīne sibbe ic forgiefe ēow.
% mdparagraph:P065:end

% mdparagraph:P066:begin
Both \textbf{ēow} forms are plural dative recipients. \textbf{Sibbe} is feminine accusative singular, the peace left or given; \textbf{mīne} agrees with it. The two first-singular verbs preserve the distinction between leaving peace and giving it. The initial objects emphasize peace, but another clear order is possible. No preposition for “with” is wanted: Cook explicitly omits it, and the taught forlǣtan frame already takes the recipient in the dative. Using forgiefan in the second clause meets the lexical control.
% mdparagraph:P066:end

\end{minipage}\par\addvspace{4pt}

\par\begin{minipage}{\linewidth}\raggedright
% mdparagraph:P067:begin
\textbf{VIII.26.5.}
% mdparagraph:P067:end

% item: EX-VIII-26-05
% key: KEY-EX-VIII-26-05
% mdparagraph:P068:begin
Ðā hit rīnde, and ðǣr cōmon flōdas, and blēowon windas, and āhruron on ðæt hūs; and ðæt hūs fēoll.
% mdparagraph:P068:end

% mdparagraph:P069:begin
The rain expression is supplied. \textbf{Flōdas} and \textbf{windas} are plural subjects of their respective past verbs and together bear on the house with \textbf{āhruron}. That verb has plural past stem -hrur-. The singular house then \textbf{fēoll}, from feallan. The on phrase expresses the goal of the assault and has accusative \textbf{ðæt hūs}. Keep explicit \textbf{ðǣr}, the narrative past and the two different fall verbs required by the prompt. Hūs is neuter in both occurrences.
% mdparagraph:P069:end

\end{minipage}\par\addvspace{4pt}

\par\begin{minipage}{\linewidth}\raggedright
% mdparagraph:P070:begin
\textbf{VIII.26.6.}
% mdparagraph:P070:end

% item: EX-VIII-26-06
% key: KEY-EX-VIII-26-06
% mdparagraph:P071:begin
Swā ǣlc gōd trēo birð gōde wæstmas, and ǣlc yfel trēo birð yfele wæstmas.
% mdparagraph:P071:end

% mdparagraph:P072:begin
Each \textbf{trēo} is singular, with third-singular \textbf{birð}. Each \textbf{wæstmas} is a masculine accusative plural object, hence \textbf{gōde} and \textbf{yfele}. The subject adjectives are neuter nominative singular \textbf{gōd, yfel}. The parallel clauses must retain both the good and the evil tree and their corresponding fruit. C01 corrects the printed prompt's wǣstm to evidenced \textbf{wæstm}; the original prompt remains available. An acceptable established beran variant may convey the same sentence, but plural berað would contradict the singular “each tree”.
% mdparagraph:P072:end

\end{minipage}\par\addvspace{4pt}

\par\begin{minipage}{\linewidth}\raggedright
% mdparagraph:P073:begin
\textbf{VIII.26.7.}
% mdparagraph:P073:end

% item: EX-VIII-26-07
% key: KEY-EX-VIII-26-07
% mdparagraph:P074:begin
Ǣlc trēo ðe gōdne wæstm ne birð forcorfen and on fȳr worpen bið.
% mdparagraph:P074:end

% mdparagraph:P075:begin
Bare \textbf{ðe} introduces the relative clause; \textbf{gōdne wæstm} is its singular masculine accusative object. The tree is the subject of a passive with two coordinated participles, \textbf{forcorfen} and \textbf{worpen}. \textbf{On fȳr} is an accusative goal without an article, as requested. Final \textbf{bið} satisfies Cook's verb-last control. [6]
% mdparagraph:P075:end

% mdparagraph:P076:begin
The established participle geworpen is a defensible alternative to worpen. Ordinary auxiliary-first passive order can be grammatical and retain the intended meaning, but it misses the stated final-verb target. A plural fruit object or an inflected relative plus ðe likewise misses a separate explicit control.
% mdparagraph:P076:end

\end{minipage}\par\addvspace{4pt}

\par\begin{minipage}{\linewidth}\raggedright
% mdparagraph:P077:begin
\textbf{VIII.26.8.}
% mdparagraph:P077:end

% item: EX-VIII-26-08
% key: KEY-EX-VIII-26-08
% mdparagraph:P078:begin
Hē ēode on munt, ðæt hē hine gebǣde; and ðā hē hine gebæd, wæs his rēaf hwīt, scīnende.
% mdparagraph:P078:end

% mdparagraph:P079:begin
\textbf{On munt} is the accusative goal of ēode. The intended prayer has past subjunctive \textbf{gebǣde}; the subsequent factual when-clause has past indicative \textbf{gebæd}. Both occurrences of gebiddan use accusative reflexive \textbf{hine}, even though English supplies no matching object. The white, shining raiment is singular neuter \textbf{rēaf}, with predicate \textbf{hwīt} and the supplied scīnende description. “He was praying” must not prompt an unlicensed literal progressive auxiliary; the clause already expresses the prayer with a finite verb. Omitting either reflexive or confusing the two moods loses a taught distinction.
% mdparagraph:P079:end

\end{minipage}\par\addvspace{4pt}

\par\begin{minipage}{\linewidth}\raggedright
% mdparagraph:P080:begin
\textbf{VIII.26.9.}
% mdparagraph:P080:end

% item: EX-VIII-26-09
% key: KEY-EX-VIII-26-09
% mdparagraph:P081:begin
Se munt, hrēosende, wierð tō nāhte.
% mdparagraph:P081:end

% mdparagraph:P082:begin
The falling mountain comes to nothing. \textbf{Wierð} is the selected third-singular present of weorðan. The participle hrēosende and the result expression \textbf{tō nāhte} are supplied for this use; neither is an invitation to invent a new independent paradigm. The participle describes the mountain without supplying a second finite predicate. Translating only “the mountain falls” omits the outcome.
% mdparagraph:P082:end

\end{minipage}\par\addvspace{4pt}

\par\begin{minipage}{\linewidth}\raggedright
% mdparagraph:P083:begin
\textbf{VIII.26.10.}
% mdparagraph:P083:end

% item: EX-VIII-26-10
% key: KEY-EX-VIII-26-10
% mdparagraph:P084:begin
Se cyning on eallum his wuldre ne wæs geglęnged swā-swā ān ðissa.
% mdparagraph:P084:end

% mdparagraph:P085:begin
The king was not adorned like one of these, despite all his glory. \textbf{On eallum his wuldre} has dative singular wuldre and matching eallum. \textbf{Ne wæs geglęnged} is a negative past passive. Cook expressly requests \textbf{swā-swā}, and \textbf{ān ðissa} means one of these, with genitive plural ðissa dependent on ān. An ordinary single swā might express comparison elsewhere but fails this local connector control. \textbf{Wuldre}, not an unchanged wuldor, is the selected syncopated dative.
% mdparagraph:P085:end

\end{minipage}\par\addvspace{4pt}

\par\begin{minipage}{\linewidth}\raggedright
% mdparagraph:P086:begin
\section*{Cook VIII.27}
% mdparagraph:P086:end
% mdparagraph:P087:begin
\textbf{VIII.27.1.}
% mdparagraph:P087:end

% item: EX-VIII-27-01
% key: KEY-EX-VIII-27-01
% mdparagraph:P088:begin
Hē ætīewde him his dēorwierðu ðing, ðæt seolfor and ðæt gold and ðā wyrtgemang and eall ðæt on his goldhordum gemētt wæs.
% mdparagraph:P088:end

% mdparagraph:P089:begin
He showed the recipients the precious things, including all four listed groups. \textbf{Him} is dative plural; the things shown are accusative. \textbf{Ætīewde} belongs to weak-I ætīewan, the verb required here, and is distinct from the earlier weak-II ætēowian. \textbf{His dēorwierðu ðing} uses the allowed strong neuter plural adjective after a possessive; a possessive does not itself make weak inflection compulsory. \textbf{His dēorwierðan ðing} is a defensible weak alternative under the earlier possessive rule. [5]
% mdparagraph:P089:end

% mdparagraph:P090:begin
\textbf{Gemētt wæs} describes what was found; \textbf{on his goldhordum} is its dative location. Do not turn the treasures into the subject of “was found”, or omit the all-that clause after listing silver, gold and spices. The possessive relations and all the listed objects matter more than copying this model's exact order.
% mdparagraph:P090:end

\end{minipage}\par\addvspace{4pt}

\par\begin{minipage}{\linewidth}\raggedright
% mdparagraph:P091:begin
\textbf{VIII.27.2.}
% mdparagraph:P091:end

% item: EX-VIII-27-02
% key: KEY-EX-VIII-27-02
% mdparagraph:P092:begin
Sum welig mann mid gōdwębbe geglęnged wæs.
% mdparagraph:P092:end

% mdparagraph:P093:begin
The subject is a certain rich man. \textbf{Mid gōdwębbe} uses the supplied material sense of mid with dative singular; the fine or purple cloth is what adorns him. \textbf{Geglęnged wæs} is past passive, with finite \textbf{wæs} last as Cook requires. \textbf{Sum welig mann wæs mid gōdwębbe geglęnged} is ordinary grammatical order, but fails that specific target. The quantity in gōdwębb is part of the verified selected spelling. [2, 6]
% mdparagraph:P093:end

\end{minipage}\par\addvspace{4pt}

\par\begin{minipage}{\linewidth}\raggedright
% mdparagraph:P094:begin
\textbf{VIII.27.3.}
% mdparagraph:P094:end

% item: EX-VIII-27-03
% key: KEY-EX-VIII-27-03
% mdparagraph:P095:begin
Ðā læg sum wǣdla æt his geate.
% mdparagraph:P095:end

% mdparagraph:P096:begin
\textbf{Wǣdla} is a singular nominative subject despite its -a ending; \textbf{læg} is singular past. \textbf{Æt his geate} has dative singular for location, and his preserves the owner of the gate. A different clear narrative order is possible. Confusing weak wǣdla with a plural would give the wrong subject and verb.
% mdparagraph:P096:end

\end{minipage}\par\addvspace{4pt}

\par\begin{minipage}{\linewidth}\raggedright
% mdparagraph:P097:begin
\textbf{VIII.27.4.}
% mdparagraph:P097:end

% item: EX-VIII-27-04
% key: KEY-EX-VIII-27-04
% mdparagraph:P098:begin
Nān mann him ne geaf; ac hundas cōmon and liccodon his wunda.
% mdparagraph:P098:end

% mdparagraph:P099:begin
No one gave to him; dogs nevertheless came and licked his sores. The negative expression is supplied, with dative \textbf{him} and an unexpressed thing given. \textbf{Geaf} uses giefan, the verb named in the omission construction taught in §5. \textbf{Hundas} is the plural of the dog noun; both \textbf{cōmon} and \textbf{liccodon} are plural past. \textbf{Wunda} is feminine accusative plural, not a masculine form inferred from the original faulty §43 pointer. C04 supplies the evidenced correction while preserving Cook's hint. [6]
% mdparagraph:P099:end

% mdparagraph:P100:begin
The thing given need not be invented. Adding a plausible “food” would add material absent from this prompt. The possessive his belongs to the poor man, while the dogs are the subjects doing the licking.
% mdparagraph:P100:end

\end{minipage}\par\addvspace{4pt}

\par\begin{minipage}{\linewidth}\raggedright
% mdparagraph:P101:begin
\textbf{VIII.27.5.}
% mdparagraph:P101:end

% item: EX-VIII-27-05
% key: KEY-EX-VIII-27-05
% mdparagraph:P102:begin
Ðā gelamp ðæt se wǣdla gewāt; and ęnglas bǣron hine tō Abrahames wīce.
% mdparagraph:P102:end

% mdparagraph:P103:begin
The supplied \textbf{gelamp ðæt} introduces what happened. \textbf{Gewāt} is the poor man's past departure; \textbf{ęnglas} are the plural bearers, and \textbf{hine} is their accusative object. \textbf{Tō … wīce} has dative singular, with the supplied genitive \textbf{Abrahames} identifying the dwelling's owner. Do not make the poor man carry the angels or change the internal possessive to a dative. Abrahames is supported here without assigning a full foreign-name paradigm.
% mdparagraph:P103:end

\end{minipage}\par\addvspace{4pt}

\par\begin{minipage}{\linewidth}\raggedright
% mdparagraph:P104:begin
\textbf{VIII.27.6.}
% mdparagraph:P104:end

% item: EX-VIII-27-06
% key: KEY-EX-VIII-27-06
% mdparagraph:P105:begin
Ðā āhōf hē his ēagan ūp and geseah Abraham feorran.
% mdparagraph:P105:end

% mdparagraph:P106:begin
\textbf{Āhōf} is the singular past of the specified āhębban; \textbf{ēagan} is the weak-neuter accusative plural object. He then \textbf{geseah} Abraham. The supplied \textbf{feorran} already expresses “from afar”; no extra preposition is required. Separate \textbf{ūp} makes the lifting-up explicit. Its omission may retain the sense in āhōf, but an ordinary hębban form would miss the prefixed lexical control. Do not replace ēagan with a regular strong-neuter plural.
% mdparagraph:P106:end

\end{minipage}\par\addvspace{4pt}

\par\begin{minipage}{\linewidth}\raggedright
% mdparagraph:P107:begin
\textbf{VIII.27.7.}
% mdparagraph:P107:end

% item: EX-VIII-27-07
% key: KEY-EX-VIII-27-07
% mdparagraph:P108:begin
Ēalā fæder Abraham, gemiltsa mē, and sęnd Lazarum.
% mdparagraph:P108:end

% mdparagraph:P109:begin
\textbf{Ēalā fæder Abraham} is direct address. The two commands are singular: have mercy on me, and send Lazarus. \textbf{Gemiltsa} takes the person in the dative, here mē; \textbf{sęnd} takes accusative \textbf{Lazarum}, supplied for this name. Mē happens to have the same surface form in dative and accusative, but the government still matters when another noun or pronoun is substituted. A nominative Lazarus does not meet the supplied accusative use.
% mdparagraph:P109:end

\end{minipage}\par\addvspace{4pt}

\par\begin{minipage}{\linewidth}\raggedright
% mdparagraph:P110:begin
\textbf{VIII.27.8.}
% mdparagraph:P110:end

% item: EX-VIII-27-08
% key: KEY-EX-VIII-27-08
% mdparagraph:P111:begin
Ēalā sunu, gemun ðæt ðū onfēnge gōdum on ðīnum līfe.
% mdparagraph:P111:end

% mdparagraph:P112:begin
The supplied \textbf{gemun} asks the son to remember an asserted past fact. \textbf{Onfēnge} is second singular past indicative, with stem onfēng- and the strong second-singular ending -e. \textbf{Gōdum} is the lexical noun “goods”, dative plural governed by onfōn. \textbf{On ðīnum līfe} is a dative singular phrase, with matching possessive. Do not parse gōdum as an accusative merely because modern English calls goods a direct object.
% mdparagraph:P112:end

\end{minipage}\par\addvspace{4pt}

\par\begin{minipage}{\linewidth}\raggedright
% mdparagraph:P113:begin
\textbf{VIII.27.9.}
% mdparagraph:P113:end

% item: EX-VIII-27-09
% key: KEY-EX-VIII-27-09
% mdparagraph:P114:begin
Ic hālsie ðē ðæt ðū hine tō mīnes fæder hūse sęnde.
% mdparagraph:P114:end

% mdparagraph:P115:begin
This uses the earlier hālsian request frame: \textbf{ðē} is the person entreated, \textbf{hine} the person to be sent. \textbf{Sęnde} is second-singular present subjunctive in the request clause. \textbf{Tō … hūse} governs dative singular house, while \textbf{mīnes fæder} remains the internal genitive phrase “of my father”. The outer preposition does not impose dative on every word inside its noun phrase. A different evidenced request verb may be valid, but merely supplying a citation form of biddan would not explain the required construction. [5]
% mdparagraph:P115:end

\end{minipage}\par\addvspace{4pt}

\par\begin{minipage}{\linewidth}\raggedright
% mdparagraph:P116:begin
\textbf{VIII.27.10.}
% mdparagraph:P116:end

% item: EX-VIII-27-10
% key: KEY-EX-VIII-27-10
% mdparagraph:P117:begin
Gif hīe ðā wītgan ne hīerað, ne gelīefað hīe, ðēah hwā of dēaðe ārīse.
% mdparagraph:P117:end

% mdparagraph:P118:begin
The plural people do not hear the prophets and will not believe. The future-context \textbf{gelīefað} needs no added “will” auxiliary. Both main propositions have their own \textbf{ne}. \textbf{Ðā wītgan} is accusative plural. The supplied indefinite \textbf{hwā} is singular subject of concessive subjunctive \textbf{ārīse}, placed last in that clause as requested. \textbf{Of dēaðe} has dative singular. An indicative would not meet this concession's taught mood; moving ārīse before its of phrase would fail the explicit word-order control even where the resulting order was otherwise defensible.
% mdparagraph:P118:end

\end{minipage}\par\addvspace{4pt}

\par\begin{minipage}{\linewidth}\raggedright
% mdparagraph:P119:begin
\textbf{VIII.27.11.}
% mdparagraph:P119:end

% item: EX-VIII-27-11
% key: KEY-EX-VIII-27-11
% mdparagraph:P120:begin
Tiliað ēow frēonda on Godes ðearfum. Gē habbað manige Lazaras nū licgende æt ēowrum gatum.
% mdparagraph:P120:end

% mdparagraph:P121:begin
The plural addressees are to obtain friends for themselves among God's poor. \textbf{Tiliað} is the plural command; \textbf{ēow} is its dative reflexive beneficiary, while \textbf{frēonda} is the genitive thing sought or obtained. \textbf{On Godes ðearfum} uses the supplied dative association/\allowbreak{}location expression. Neither “for yourselves” nor “friends” should be mechanically turned into an English-style accusative object. [5]
% mdparagraph:P121:end

% mdparagraph:P122:begin
The second sentence says that many people like Lazarus are now lying at their gates. \textbf{Manige Lazaras} is the accusative plural object of habbað, with supplied \textbf{licgende} describing that object. It is not a new compound tense formed from habban and a participle. \textbf{Æt ēowrum gatum} has dative plural agreement and the selected short-stem plural gatum. The proper-name plural is supplied; its appearance does not assign an independent Latin-name declension.
% mdparagraph:P122:end

\end{minipage}\par\addvspace{4pt}

\clearpage
% mdparagraph:P123:begin
\section*{S. The Dangers of Greatness}
% mdparagraph:P123:end

\reserve{5\baselineskip}
\reserve{8\baselineskip}
% mdparagraph:P124:begin
\textbf{S01. Complete translation.}
% mdparagraph:P124:end

% item: NEW-U14-S01
% key: KEY-NEW-U14-S01
% mdparagraph:P125:begin
Every high rank here in the world is surrounded by dangers; the greater the worldly honours, the fiercer the dangers. Thus, from that, we can recognize and understand the examples.
% mdparagraph:P125:end

% mdparagraph:P126:begin
The tree, then, which grows in the woodland so that it towers above all the other trees and spreads itself out—when a storm and the strong wind suddenly come on, it is shaken and battered more violently than the rest of the forest.
% mdparagraph:P126:end

% mdparagraph:P127:begin
It is likewise with the high cliffs and crags: when they tower far above the rest of the earth, they then suddenly begin to fall and crash very violently to the ground.
% mdparagraph:P127:end

% mdparagraph:P128:begin
Likewise with the high mountains and hills, which stand high above the whole world: nevertheless they suffer the penalty of their eminence, in that they are often chastised and afflicted by fire from heaven and struck by thunderbolts.
% mdparagraph:P128:end

% mdparagraph:P129:begin
So the high powers here in the world collapse and fall and come to ruin, and the riches of this world turn to sorrow, and these earthly wonders come to nothing.
% mdparagraph:P129:end

% mdparagraph:P130:begin
“High honour”, “office” or “position” can express ār here; “worldly dignities/\allowbreak{}distinctions” can express woruldgeðyngða. Preserve the proportional claim: greater rank brings greater danger. \textbf{Be ðǣm} refers separately to the preceding claim; \textbf{ðā bȳsena} is the accusative object of both “recognize” and “understand”. “From these examples” merges two distinct phrases and misses that syntax. [7]
% mdparagraph:P130:end

% mdparagraph:P131:begin
“Groves” or “woodland grove” can render wudubearwe. The storm and strong wind are coordinated subjects of the source's singular gestęnt; English can use a plural verb. “Set in”, “blow up” or “strike” can convey the storm's onset/\allowbreak{}assault in context. The tree spreads itself and then receives the shaking. “The rest of the forest” is not a different forest. “Crags” retains Cook's probable contextual interpretation of torrum; the following “rest of the earth” has its own supplied explanation despite the original note's misplaced reference. [7]
% mdparagraph:P131:end

% mdparagraph:P132:begin
The cliff infinitives can be rendered “fall and collapse” or “begin to fall and crash”; they reinforce one downfall. “Pre-eminence”, “supremacy” or “chief place” expresses the mountains' ealdordōm in this comparison. The passive clause explains their penalty; it does not state their purpose. \textbf{Līegum} means thunderbolts/\allowbreak{}lightning in the selected glossary. Rendering both heofonfȳre and līegum simply as “flames” loses the latter lexical distinction. “Come to naught” is a natural alternative to “come to nothing”. Keep all three human outcomes: powers to ruin, riches to sorrow, wonders to nothing. [7]
% mdparagraph:P132:end

\reserve{5\baselineskip}
\par\begin{minipage}{\linewidth}\raggedright
% mdparagraph:P133:begin
\textbf{S02. The argument.}
% mdparagraph:P133:end

% item: NEW-U14-S02
% key: KEY-NEW-U14-S02
% mdparagraph:P134:begin
The opening claim is that high worldly rank brings danger, and increased rank brings greater danger. The three examples proceed from the tall spreading tree, through high cliffs/\allowbreak{}crags, to high mountains/\allowbreak{}hills. The exposed tree suffers stronger shaking; cliffs fall violently to earth; mountains pay for their height by suffering heavenly fire and thunderbolts. The conclusion applies these images to human powers, worldly wealth and earthly wonders, whose height or splendour does not prevent their downfall.
% mdparagraph:P134:end

% mdparagraph:P135:begin
A concise paraphrase is sufficient if it preserves the order and connection. Listing three natural objects without explaining their relation to rank and danger misses the argument; adding the following unassigned luxury/\allowbreak{}death paragraph is unnecessary.
% mdparagraph:P135:end

\end{minipage}\par\addvspace{4pt}

\par\begin{minipage}{\linewidth}\raggedright
% mdparagraph:P136:begin
\textbf{S03. The tree sentence.}
% mdparagraph:P136:end

% item: NEW-U14-S03
% key: KEY-NEW-U14-S03
% mdparagraph:P137:begin
All three occurrences of \textbf{hit} refer to the tree: the first is the nominative subject of \textbf{hlīfað}; the second is the reflexive accusative object of \textbf{brǣt}, “spreads itself”; the third is the nominative subject of \textbf{bið … gewǣged and geswęnged}, “is shaken and battered”. Identical spelling does not imply an identical case or role.
% mdparagraph:P137:end

% mdparagraph:P138:begin
\textbf{Ðe wiexð on ðǣm wudubearwe} identifies the tree that grows in the woodland. \textbf{Ðæt hit hlīfað … and brǣt hit} gives the result of its growth: it towers above the other trees and spreads out. The initial parenthetic ðonne after trēow is “then” in the exposition. The later \textbf{ðonne … ðonne} is temporal and correlative: “when the storm and wind come on … then the tree is shaken”. English can omit the resumptive “then”. The final \textbf{ðonne} means “than” and introduces the comparison with \textbf{se ōðer wudu}, “the rest of the forest”. Do not translate all these occurrences with the same English connective. [7]
% mdparagraph:P138:end

\end{minipage}\par\addvspace{4pt}

\par\begin{minipage}{\linewidth}\raggedright
% mdparagraph:P139:begin
\textbf{S04. The cliffs.}
% mdparagraph:P139:end

% item: NEW-U14-S04
% key: KEY-NEW-U14-S04
% mdparagraph:P140:begin
Both \textbf{hīe} refer to the high cliffs and crags. \textbf{Onginnað} governs both infinitives, \textbf{feallan} and \textbf{hrēosan}. The two near-synonymous verbs emphasize their fall or collapse; they need not describe two successive falls. \textbf{Full ðearlice}, “very violently”, intensifies that description, particularly hrēosan. A natural rendering is “they suddenly begin to fall and crash very violently to the ground”. Do not attach hrēosan to a different unexpressed finite verb. Cook's historical note is reference material; no knowledge of Chaucer or Greek is required. [7]
% mdparagraph:P140:end

\end{minipage}\par\addvspace{4pt}

\par\begin{minipage}{\linewidth}\raggedright
% mdparagraph:P141:begin
\textbf{S05. Source derivatives.}
% mdparagraph:P141:end

% item: NEW-U14-S05
% key: KEY-NEW-U14-S05
% mdparagraph:P142:begin
\textbf{Ealdordōmes:} \textbf{ealdor + -dōm + -es}. The derivational element -dōm forms the noun denoting leadership or primacy; -es is genitive singular inflection. Here the noun denotes the mountains' eminence. \textbf{Wīte … ðæs ealdordōmes} is the penalty of that pre-eminence: their height brings exposure to the heavenly assaults. A contextual “penalty for their supremacy” is acceptable; ealdordōmes is not the subject of habbað.
% mdparagraph:P142:end

% mdparagraph:P143:begin
\textbf{Eorðlican:} \textbf{eorð- + -lic + -an}. The derivational -lic forms the adjective “earthly” from the earth family; -an is its weak nominative neuter plural ending here, agreeing with \textbf{ðās … wundor}, “these earthly wonders”. Its clause determines nominative subject function. The same -an spelling elsewhere would not by itself fix case, number or gender. Neither analysis requires producing unassigned forms of either word. [4, 7]
% mdparagraph:P143:end

\end{minipage}\par\addvspace{4pt}

\par\begin{minipage}{\linewidth}\raggedright
% mdparagraph:P144:begin
\textbf{S06. The conclusion's three subjects.}
% mdparagraph:P144:end

% item: NEW-U14-S06
% key: KEY-NEW-U14-S06
% mdtable:T12:begin
\par\begingroup
\fontsize{10.5}{13.2}\selectfont
\setlength{\LTleft}{0pt}\setlength{\LTright}{\fill}
\noindent
\begin{tabular}{@{}>{\raggedright\arraybackslash}p{30.05mm}>{\raggedright\arraybackslash}p{91.13mm}@{}}
\toprule
\textbf{Subject} & \textbf{Outcome} \\
\midrule
ðā hēan mihta & hrēosað, feallað, tō lore weorðað: the high powers collapse, fall and come to ruin. \\
ðisse worulde welan & weorðað tō sorge: the riches of this world turn to sorrow. \\
ðās eorðlican wundor & weorðað tō nāhte: these earthly wonders come to nothing. \\
\bottomrule
\end{tabular}
\par\endgroup

% mdtable:T12:end
% mdparagraph:P145:begin
\textbf{Ðisse worulde} is a dependent feminine genitive singular phrase, “of this world”, within the second subject. Its head is \textbf{welan}, the plural subject “riches”. \textbf{Ðās} directly determines \textbf{wundor} and agrees with it in nominative neuter plural. A plural subject may contain a singular genitive dependent; the dependent phrase does not set the number of weorðað. Keep the riches and wonders as subjects of their own clauses, not objects of the preceding falling verbs. [7]
% mdparagraph:P145:end

\end{minipage}\par\addvspace{4pt}

\par\begin{minipage}{\linewidth}\raggedright
% mdparagraph:P146:begin
\textbf{S07. Pairs and repetitions.}
% mdparagraph:P146:end

% item: NEW-U14-S07
% key: KEY-NEW-U14-S07
% mdparagraph:P147:begin
Any two actual examples with a defensible contextual explanation suffice. \textbf{Oncnāwan and ongietan} reinforces recognizing and understanding the examples. \textbf{Feallan … hrēosan} in the cliff sentence, or \textbf{hrēosað, and feallað} in the conclusion, reinforces collapse without requiring two separate events. The three uses of \textbf{weorðað} with \textbf{tō lore, tō sorge, tō nāhte} align the ruin, sorrow and nothingness of worldly possessions and power; the order of verb and tō phrase varies. Repeated height vocabulary, such as \textbf{hēah ār}, \textbf{hēaclifum}, \textbf{hēagum muntum} and \textbf{hēan mihta}, joins social rank to the natural illustrations.
% mdparagraph:P147:end

% mdparagraph:P148:begin
Another genuine pair or repetition is acceptable if quoted accurately and tied to this argument. Naming a general device without quoting the passage, or asserting that every repeated verb describes an additional event, misses the task.
% mdparagraph:P148:end

\end{minipage}\par\addvspace{4pt}

\clearpage
% mdparagraph:P149:begin
\section*{Sources and selected forms}
% mdparagraph:P149:end

\begingroup\fontsize{10}{12.8}\selectfont
% mdparagraph:P150:begin
\textbf{1. Scope and assessment.} Cook's \emph{Exercises in Old English} (1895), Group VIII, Exercises 26–27, printed pp. 13–15 /\allowbreak{} one-based PDF 21–23, supplies the unchanged English prompts. The Old English responses and the original exercise answers are new teaching compositions. The unit's separately defined morphology and lexical frames govern its selected models. The source and supplied helps remain available during supported work.
% mdparagraph:P150:end

% mdparagraph:P151:begin
\textbf{2. Nouns and reference cells.} Cook, \emph{A First Book in Old English}, second edition (1895), §§43, 47, 51 and 53, printed pp. 28–37 /\allowbreak{} PDF 44–53. Section 53.2 gives the weak-neuter ēage relation and alternate genitive plural ēagna. Section 51.5 explains -nes with oblique -ness-. Selected lexical quantities are checked against the glossary: dūst at p. 283 /\allowbreak{} PDF 299, gōd-wębb at p. 296 /\allowbreak{} PDF 312. Tables mark abstract/\allowbreak{}material number references separately from ordinary production.
% mdparagraph:P151:end

% mdparagraph:P152:begin
\textbf{3. Verbal systems.} Cook §§93–95, pp. 54–55; strong-II/\allowbreak{}III models pp. 60–62; strong V–VII pp. 63–67; weak I pp. 68–72; weak II §§118–119, pp. 73–74; dental changes pp. 24–25. These rules and the individual lexical entries support the selected complete systems. Past participle prefixes and lexical ge- remain distinct. Bosworth–Toller, main p. 303b, for-ceorfan, corroborates forceorfende and forcorfen. Citation diagnostics do not assign unlisted syntactic uses.
% mdparagraph:P152:end

% mdparagraph:P153:begin
\textbf{4. Word formation.} Cook §§142–147, pp. 81–87, supplies the bounded prefix, suffix and compound analysis; §70, p. 44, supplies the selected adverb families. The individual glossary entries establish whole-word meanings. The source ealdordōmes and eorðlican have their supplied contextual analyses; no additional productive paradigms follow from their segmentation.
% mdparagraph:P153:end

% mdparagraph:P154:begin
\textbf{5. Government and earlier scope.} Cook's dative and genitive constructions, pp. 91–95, and prepositions, pp. 96–99, are read with the actual lexical entries. The selected onfōn and gemiltsian dative examples occur on p. 94. Bosworth–Toller, main p. 411b, ge-lýfan, directly corroborates belief on + accusative in on hine gelýfan; the course retains Cook’s selected ie spelling. Reflexive uses under §184.a/\allowbreak{}b are on pp. 100–101. Earlier Unit 11 explicitly teaches the article-plus-weak rihtwīs personal expression; Unit 12 teaches the sęndan recipient and command frames; Unit 13 teaches the hālsian request. Unit 3 allows strong or weak adjectives after a possessive; weak inflection is not compulsory solely because his precedes the adjective.
% mdparagraph:P154:end

% mdparagraph:P155:begin
\textbf{6. Original controls and corrections.} Cook §87.d, p. 52 /\allowbreak{} PDF 68, gives bare relative ðe. Purpose under §196.f, p. 105 /\allowbreak{} PDF 121, is distinct from asserted result. Preserve the final-verb controls in 26.7, 27.2 and the concessive clause of 27.10. Correction C01 retains original wǣstm in 26.6 while selecting evidenced wæstm. C04 retains 27.4's original wund, 43 but identifies the feminine wound noun: Bosworth–Toller main p. 1281a and Supplement p. 570b, where the relevant Luke 16.21 citation has his wunda liccodon. Cook's §43 is masculine; §51.b, p. 35, supplies the selected feminine pattern. The adjective wund in Cook's glossary is a distinct word.
% mdparagraph:P155:end

% mdparagraph:P156:begin
\textbf{7. Reading.} Cook, Reader VIII, \emph{The Dangers of Greatness}, printed 156.1–158.4 /\allowbreak{} PDF 172–174, is the complete assigned opening argument, 159 words. Prose-line counts are editorial, anchored by printed marginal numbers. The following luxury/\allowbreak{}death argument is unassigned. All 24 attached original notes and 25 anchors are retained in the learner chapter. Source note 156.9 identifies wind as the second subject of gestęnt; 156.11 explains se ōðer wudu as the rest of the forest. C05 preserves the apparent mismatch in 157.1's printed pointer to 156.12 and identifies 156.11 as the relevant parallel. Cook p. 303 /\allowbreak{} PDF 319 glosses līeg as thunderbolt/\allowbreak{}levin. The source's qualified attribution is retained; its literary parallels are Cook's suggestions, not newly established borrowings.
% mdparagraph:P156:end

\endgroup
\end{document}
